\section{Single insertions, quadratic charges, and time cones}
\label{sec:charges-cones}

Throughout this Section, $l$ is the dimension-two vacuum wave constructed in
Section~\ref{sec:foundations}. The letter $L$ will initially denote only an
insertion in a sesquilinear Wightman form. Neither an operator $L(f)$ nor a
product containing two copies of $L$ is being asserted. This distinction is
essential when conserved currents are used below.

The section proceeds through three distinct constructions.  We first
extend the extreme vacuum values of $l$ to local forms with one
insertion.  Their conserved currents then define charge forms, whose
quadratic part is identified with the positive moment $K_n$ on a
specified domain.  Finally the vanishing residual dilation charge
identifies the actual cone modular groups and produces the actual
time-M\"obius action used for scalar extraction.

\subsection{From scalar localization to one local insertion}

The order-zero spectral comparison has a localization interpretation.
There are two transports: from the scalar comparison to the moment
$\kappa$, and from that moment to a physical bounded cut. We state both,
together with the reverse containment for actual parabolic cuts that
will be needed later.

\begin{lemma}[Localization transport]\label{C:localization-transport}
Let $X_0(W)$ be the scalar wedge real space in the $l$ channel,
$\iota$ its embedding into the physical Hilbert space, and $X(W)$ the
physical wedge real space. For $p(z)=c|z-d|^2+c_0$, $c>0$, put
\[
 K_{0,p}=b\{c(|Y-d|^2+h_0)+c_0\},\qquad
 K_p=b\{c(|Y-d|^2+\kappa)+c_0\}.
\]
Then
\begin{equation}\label{C:moment-localization}
 \iota e^{iK_{0,p}}X_0(W)\subset e^{iK_p}X(W).
\end{equation}
If $g$ is bounded and smooth and $g\le p$, the left side also lies in
$X(W_g)=e^{iH_g}X(W)$. Conversely, the actual parabolic real space
satisfies
\begin{equation}\label{C:actual-parabolic-localization}
 X(W_p)\subset e^{iK_p}X(W).
\end{equation}
These statements hold after translations of the null plane and reflection.
\end{lemma}

\begin{proof}
The scalar wedge conjugation is the real scalar conjugation; it is
identified already on the dense waves generated by $M^2l$. To prove
\eqref{C:moment-localization}, cancel the common external factors and
consider $e^{-icb\kappa}\iota e^{icbh_0}$. Let $E_i$ be the
actual $\kappa$ step projections and $P_j$ the reference $h_0$ step
projections. Clip both spectral variables at one common finite bound,
then approximate the clipped
$\kappa$ from below and the clipped $h_0$ from above.  If their step
values are $a_i,c_j$, spectral inclusion gives
$E_i\iota P_j=0$ whenever $a_i>c_j$.  Each surviving ordered term is
therefore $E_i\iota P_j e^{icb(c_j-a_i)}$, with nonnegative
translation parameter; no joint spectral resolution is used.
The coefficients intertwine the boost and wedge conjugation, so
they map the reference real wedge space into the actual one.  The
positive translations preserve that actual space.  At the common
clipping endpoint both step values coincide, which preserves the
same sign.  Strong convergence follows from the separate spectral
calculi as the mesh tends to zero and the clipping bound tends to
infinity. Closedness of $X(W)$ proves the inclusion. This uses the
ordered products of two separate spectral resolutions, not a joint
resolution of $\kappa$ and $h_0$.

For the second transport, positivity, normalization, and the zero first
moment of $\mu$ give
\[
 H_g=b\int g(Y+z)\mu(\dd z)
 \le b\{c(|Y-d|^2+\kappa)+c_0\}=K_p.
\]
This inequality is in the common commuting calculus. Its difference
is $b$ times a nonnegative boost-independent operator commuting with
wedge conjugation. Thus $e^{i(K_p-H_g)}$ is a direct integral of
inward null translations and preserves $X(W)$. Apply it to
\eqref{C:moment-localization} to obtain membership in $X(W_g)$.

Finally the actual parabolic generator is
$H_p=b\{c(|Y-d|^2+h)+c_0\}$. Since $h$ and $\kappa$ commute and
$\kappa\le h$, the same argument applies to
$H_p-K_p=bc(h-\kappa)\ge0$ and proves
\eqref{C:actual-parabolic-localization}. The operators vanish on
$b=0$. Covariance and reflected agreement give the remaining placements.
No equality of $h$ with $\kappa$ is used.
\end{proof}

Here is the resulting extreme-ordering statement. First let
$A_i(y_i)$ be homogeneous original physical fields. The two
expressions with $L(x)$ at the extreme left or
extreme right are defined by pairing the wave $l$ with a physical
polynomial vector. They agree whenever $x$ is spacelike to all $y_i$ and,
for some future null $n$,
\begin{equation}\label{C:extreme-patch}
 n\cdot x<\min_i n\cdot y_i
 \quad\hbox{or}\quad
 n\cdot x>\max_i n\cdot y_i.
\end{equation}
For the first alternative, place a null plane between $x$ and the
$y_i$ and use the future shadow of a tip just before $x$ as a dominating
parabola. The past shadows of the $y_i$ lie strictly below it by
spacelikeness. Their differences from it tend quadratically to
$-\infty$ at transverse infinity, so a bounded smooth truncation from
below still separates the supports. These strict inequalities persist
in small neighborhoods. Locality of bounded operations in that cut,
followed by Tomita adjunction of the whole physical product, proves the
equality. The opposite alternative uses the reflected construction.
The scalar insertion is even, so no additional graded sign occurs.
For a permitted compact joint kernel, the same statement follows by
the tensor-test approximation and closed Tomita-graph limit in the
discussion of \eqref{F:BW-sandwich}; it does not require an affiliated
realization for the joint-smeared polynomial itself.

The remaining task is to place this insertion between arbitrary
physical factors while retaining their ordered spectra. Causal solutions
at adjacent interchanges will have the required wave source, but their
compatibility leaves a homogeneous-wave defect. The next lemma recovers
such a defect from its support and spectator analyticity. We then apply
it to assemble the insertion with all its ordering cones.

\Needspace{7\baselineskip}
\begin{lemma}[Recovery of a homogeneous wave defect]
\label{C:homogeneous-defect}
Let $H(x,y_1,\ldots,y_m)$, $m\ge1$, be a tempered, jointly
translation-invariant distribution with $\Box_xH=0$ and
\begin{equation}\label{C:H-support}
 \supp_x H\subset C_y\cup K_y,\qquad
 C_y=\bigcup_iJ(y_i),\qquad
 K_y=\bigcap_{n\ {\mathrm{future\ null}}}
 \{x:\min_i n\cdot y_i\le n\cdot x\le\max_i n\cdot y_i\}.
\end{equation}
For a compact transfer-momentum test $\varphi(k)$, let $H_\varphi(y)$
be its pairing with the partial Fourier transform in $x$, using
$e^{-ikx}$. Suppose these pairings are holomorphic in the ordered
spectator tube and, for $\supp\varphi\subset\{|k|\le K\}$, $0<K\le1$,
obey
\begin{equation}\label{C:shifted-test-bound}
 |H_\varphi(y)|\le C p_N(\varphi)
 (1+|y|+\delta(y)^{-1})^N
 \exp(CK|\operatorname{Im}y|).
\end{equation}
Here $p_N$ is one fixed finite Schwartz seminorm and $\delta(y)$ is the minimum
distance of the successive imaginary gaps from the cone boundary
($\delta=1$ for $m=1$). Then $H$ is a finite sum
\[
 H(x,y)=\sum_{i,P}C_{iP}(y)P(\partial_x)\Delta_0(x-y_i),
\]
where the coefficients have tempered boundaries with the ordinary
ordered spectator spectra. They are invariant under common spectator
translation. The polynomials $P$ may be chosen independently modulo
$k^2$. We normalize the massless commutator solution by
$\Delta_0(0,\boldsymbol r)=0$ and
$\partial_0\Delta_0(0,\boldsymbol r)=\delta(\boldsymbol r)$.
\end{lemma}

\begin{proof}
Normalizing $n=(1,\omega)$ and taking $\omega=\pm e_a$ shows that
$K_y$ is compact, uniformly on compact spectator sets. We first remove
its part outside $C_y$, then recover the coefficients of the remaining
wave solutions with bounds that preserve the unshifted spectator spectra.

We give the support argument explicitly. Choose an open spectator
patch around distinct equal-time spatial extremal vertices, so that
each $y_i$ is a strict maximum and a strict minimum for open sets of
future null directions. If $x\in K_y\cap J^+(y_i)$, the maximizing
directions give $0\le n\cdot(x-y_i)\le0$. A causal vector orthogonal
to an open set of null directions is zero. The minimizing directions
treat $J^-(y_i)$. Hence
\begin{equation}\label{C:vertex-intersection}
 K_y\cap C_y=\{y_1,\ldots,y_m\}
\end{equation}
on a smaller patch. Fix a test neighborhood $U$ spacelike to all the
vertices, and choose small vertex neighborhoods $N_i$ whose closures
remain spacelike to $U$. There is a compactly supported smooth
$\chi(x)$ equal to one near $K_y$ whose derivatives meet $C_y$ only
in $N=\bigcup_iN_i$. Positive separation margins make one such cutoff
work uniformly after shrinking the spectator patch. A compact
spectator cutoff causes no change to $\Box_x$. Thus $v=\chi H$ has
compact $x$-support and
\[
 \Box_xv=[\Box_x,\chi]H=:f,\qquad \supp_xf\subset N.
\]
For compactly supported distributions the exact identity
$v=E_{\rm ret}*\Box_xv$ follows by moving $\Box$ through convolution.
Its right side vanishes on $U$, since $J^+(N)\cap U=\varnothing$.
Near $K_y$ one has $v=H$; elsewhere the original support statement
already suffices. This removes $K_y\setminus C_y$ from
\eqref{C:H-support}, without any decomposition of distributions on
touching closed sets.

In a common thin time slab the cone components are disjoint.
A spatial partition separates them with derivatives off their
supports. Each component still solves the wave equation there.
In $r=x-y_i$ coordinates its displacement and normal-derivative data
on $r^0=0$ are supported at $\boldsymbol r=0$. Distributional
solutions of the wave equation have distributional Cauchy data on
this non-characteristic slice. The point-support structure theorem
therefore writes both data as finite sums of spatial delta
derivatives with spectator-distribution coefficients. On a smaller
relatively compact patch their orders are uniformly bounded.
With $\Delta_0(0,\boldsymbol r)=0$ and
$\partial_0\Delta_0(0,\boldsymbol r)=\delta(\boldsymbol r)$,
their evolutions are derivatives of $\Delta_0$. Cauchy uniqueness
first in the slab and then globally in $x$ gives
\begin{equation}\label{C:vertex-wave-expansion}
 H(x,y)=\sum_{i,P}C_{iP}(y)P(\partial_x)\Delta_0(x-y_i).
\end{equation}
Choose the finitely many polynomials $P$ independently modulo $k^2$.
Enlarging the finite family if necessary, use homogeneous bases
$P_{q\alpha}$ of the quotient by $k^2$ in each degree
$q=0,\ldots,N$.

We next justify continuation of the coefficients, including its
growth bound. This is needed to recover the unshifted spectator
spectra. Smear the Fourier transform of
$P(\partial_x)\Delta_0(x-y_i)$ against compact $k$-tests. Its column
is, up to constant factors,
\[
 \big\langle P(k)\widehat\Delta_0(k)e^{ik y_i},\varphi(k)\big\rangle.
\]
All augmented minors expressing that the corresponding tests of
$H$ lie in the span of these columns vanish on the initial patch.
They are holomorphic in the ordered spectator tube; boundary
uniqueness makes them vanish throughout that tube. The columns have
full rank at every interior point, with a quantitative bound as
follows. Choose finitely many fixed positive null vectors
$n_r=(1,\omega_r)$ such that, for every $q\le N$, the evaluation
matrix
\[
 (E_q)_{r\alpha}=P_{q\alpha}(n_r)
\]
is injective. Such a choice exists: a homogeneous polynomial
vanishing on every positive null ray is zero modulo $k^2$, and
finite dimensionality permits a finite separating set of
evaluations. Taking the union of these sets for $q\le N$ gives
one collection of rays. Fix a left inverse for each $E_q$.

On the $r$th ray the polynomial-exponential column sum is
\[
 \sum_{i=1}^m\sum_{q=0}^N
 a_{iqr}s^qe^{is n_r\cdot y_i},
 \qquad
 a_{iqr}=\sum_\alpha C_{i,q\alpha}(y)P_{q\alpha}(n_r).
\]
Let $d=m(N+1)$, $\varepsilon=c/(1+|y|)$, and use the
$d$ samples $s=t\varepsilon$, $t=1,\ldots,d$. The matrix
from the coefficients $a_{iqr}$ to these samples is
\[
 (V_r)_{t,(i,q)}=(t\varepsilon)^qz_{ir}^{\,t},
 \qquad z_{ir}=e^{i\varepsilon n_r\cdot y_i}.
\]
Its determinant is the confluent Vandermonde determinant
\begin{equation}\label{C:Vandermonde}
 \det V_r=c_{m,N}\varepsilon^{mN(N+1)/2}
 \prod_i z_{ir}^{(N+1)(N+2)/2}
 \prod_{i<j}(z_{jr}-z_{ir})^{(N+1)^2},
 \qquad c_{m,N}=\Bigl(\prod_{q=0}^Nq!\Bigr)^m.
\end{equation}
For example this follows by replacing powers of $t$ by falling
factorials and using the derivative columns of the ordinary
Vandermonde matrix. Thus the ray inversion distinguishes the
degree $q$ and center $i$ before the fixed left inverses of
$E_q$ recover each coefficient $C_{i,q\alpha}$. This describes
a left inverse for the full rectangular multi-ray matrix;
independence on each ray alone would not suffice.

The scale of the samples is also significant.  Keeping them at fixed
nonzero momenta would preserve the exponential type of the shifted
spectator cones.  Taking $s=O((1+|y|)^{-1})$ instead keeps that
exponential bounded.  The determinant estimate ensures that this
shrinking choice costs only polynomial powers, which are compatible
with tempered spectral boundary values.

For the tube distance $\delta(y)$ in the statement, the moduli
of the $z_{ir}$ are bounded
above and below, and for $i\ne j$
\[
 |z_{ir}-z_{jr}|\ge\big||z_{ir}|-|z_{jr}|\big|
 \ge c'\varepsilon\delta(y)\quad(i\ne j),
\]
because the imaginary gaps are strictly ordered and the directions
are fixed. The determinant formula and cofactor bounds therefore
give the full left inverse polynomial norm in
$\varepsilon^{-1}$ and $1+\delta(y)^{-1}$. To turn these point
samples of the known columns into legitimate compact $k$-tests,
replace each sample $t\varepsilon n_r$ by a smooth bump of width $w$,
normalized to have integral one on the positive null shell. The centers
are a distance comparable to $\varepsilon$ from its vertex. For
$w\le c_1\varepsilon$, its shell Jacobian and the prescribed test
derivatives therefore give
\[
 p_N(\varphi)\le C\varepsilon^{-a}w^{-b}
\]
for fixed finite $a,b$. Derivatives of the column functions
$P(k)e^{ik\cdot y}$ cost polynomial powers of $1+|y|$ and
$\varepsilon^{-1}$, while their exponential is bounded on these
supports. If $A_0$ is the point-sample matrix and $L_0A_0=1$,
choose
\[
 w=c_2\varepsilon(1+|y|+\delta(y)^{-1})^{-J}
\]
with fixed $J$ sufficiently large and $c_2$ sufficiently small. The
resulting test-column matrix $A$ then satisfies
$\|L_0(A-A_0)\|<1/2$. Its left inverse
$(L_0A)^{-1}L_0$ has the same polynomial bound, and the displayed
seminorm cost remains polynomial.

This proves full rank using genuine test columns. On a neighborhood
of any fixed interior point, keep a full-rank finite set of tests
fixed; its inverse and the vanishing augmented minors define
holomorphic coefficients. Uniqueness identifies them on overlaps.
The parameter-dependent test choices above serve only to estimate
these already defined holomorphic functions. Their supports satisfy
$|k|=O(\varepsilon)$ and their seminorms have polynomial cost.
Applying \eqref{C:shifted-test-bound} with $K=O(\varepsilon)$
therefore costs only polynomial powers: the remaining exponential
is bounded by $\exp(O(\varepsilon|\operatorname{Im}y|))=O(1)$.
The left inverse consequently gives
\begin{equation}\label{C:coefficient-bound}
 |C_{iP}(y)|\le C(1+|y|)^N(1+\delta(y)^{-1})^N
\end{equation}
for some $C,N$ on the tube. The coefficients have ordinary tempered
ordered spectral boundaries. Testing all compact $k$-sets, including
a neighborhood of zero, proves \eqref{C:vertex-wave-expansion}
globally. Translation covariance also shows that the coefficients
are invariant under common spectator translation.

\end{proof}

The construction is first carried out for a fixed list of homogeneous
original entries.  We will then extend it by finite multilinearity in
their scaling components.  Descent to polynomial vectors will be
checked for the whole summed bra or ket, so cancellations between
different homogeneous terms remain allowed.

\Needspace{6\baselineskip}
\begin{proposition}[One-insertion forms]\label{C:single-insertion}
There are unique tempered, Hermitian, Poincar\'e- and scale-covariant
Wightman forms with exactly one insertion $L$ and arbitrarily many
physical insertions, whose extreme vacuum values are those of $l$ and
which satisfy $\Box L=\Theta$. They have the ordinary ordered spectra,
are local under every allowed adjacent interchange, and depend only on
the physical polynomial vectors on the two sides of $L$.
\end{proposition}

\begin{proof}
\emph{Causal source solutions.}
The empty physical product has only the specified zero vacuum
expectation, since $l$ has no zero-momentum component. Suppose
$m\ge1$ and fix an order of homogeneous original entries
$A_i\in\mathcal F_{\Delta_i}$, $i=1,\ldots,m$.
Use the Fourier convention
$e^{-ikx-i\sum_i p_i y_i}$, and put $q_h=\sum_{i\le h}p_i$, with
$q_0=0$. For an insertion just after $A_j$, the required spectral
conditions are
\begin{equation}\label{C:insertion-spectra}
 q_h\in\overline V_+\quad(1\le h\le j),
 \qquad k+q_h\in\overline V_+\quad(j\le h\le m-1),
 \qquad k+q_m=0.
\end{equation}
Redundant zero partial sums at the endpoints may be included. Write
$E_R,E_L$ for the known extreme values. For each $i$, take the
difference of the physical distributions with $\Theta$ immediately
before and immediately after $A_i$. Its support, in $r=x-y_i$, is
causal. Split it into forward and backward parts by a finite-order
Taylor subtraction in $r$ alone, and apply respectively the retarded
and advanced wave fundamental solutions. This gives a tempered $B_i$
with the same causal support and the prescribed wave source. The
convolutions are proper on the respective cones: two future causal
vectors, or two past causal vectors, are bounded on the inverse image
of a bounded set under addition. The same cone geometry gives the
polynomial bounds needed for tempered convolution.

This construction preserves precisely the common spectral cones
\begin{equation}\label{C:merged-spectra}
 q_h\in\overline V_+\ (h<i),\qquad
 k+q_h\in\overline V_+\ (h\ge i).
\end{equation}
Indeed in the coordinates $(r,y_1,\ldots,y_m)$ the spectator momenta
are $p_j$ for $j\ne i$ and $k+p_i$ for $j=i$.
All the conditions in \eqref{C:merged-spectra} involve only those
spectator momenta. Multiplication, subtraction, and convolution in $r$
do not alter them. This is why solving the wave equation separately
in the external momentum before doing causal splitting would be
insufficient.

\emph{Removing the homogeneous defect.}
Set
\[
 R_j=E_R+\sum_{i>j}B_i,\qquad
 L_j=E_L-\sum_{i\le j}B_i.
\]
Both have the wave source appropriate to the $j$th position. The first
has all the initial cones in \eqref{C:insertion-spectra}; the second
has all its final cones. Their difference
$H_*=R_j-L_j$ is independent of $j$ and satisfies $\Box_xH_*=0$.
After a compact Fourier test in $k$, its spectator partial sums lie in
$\overline V_+$ up to bounded shifts of size $O(\sup|k|)$.
It consequently has an analytic continuation to the ordinary ordered
spectator tube, with polynomial bounds multiplied by an exponential
of type $O(\sup|k|)$. Equation~\eqref{C:extreme-patch} and the causal
support of the $B_i$ give
the support condition $\supp_xH_*\subset C_y\cup K_y$ of
\eqref{C:H-support}.

Its shifted spectator cones also give
\eqref{C:shifted-test-bound}: for a transfer test supported in
$|k|\le K$, every partial cone is shifted by at most a fixed multiple
of $K$. Fourier--Laplace damping contributes
$\exp(CK|\operatorname{Im}y|)$; the finite-order tempered bound on the
joint distribution controls the remaining terms by one fixed seminorm
$p_N(\varphi)$ and polynomial powers of $|y|$ and the inverse tube
distance. This estimate is uniform in the test, including when its
support shrinks toward $k=0$. Lemma~\ref{C:homogeneous-defect} therefore
gives the expansion \eqref{C:vertex-wave-expansion} for $H_*$.

Let $H_i$ denote the summand centered at $y_i$. In Fourier variables
it shifts $p_i$ by $k$, and no other spectator momentum. Consequently
it satisfies exactly the common cones
\eqref{C:merged-spectra}. The required distributions are
\begin{equation}\label{C:Wj}
 W_j=R_j-\sum_{i>j}H_i
     =L_j+\sum_{i\le j}H_i.
\end{equation}
The first expression supplies the initial cones and the second the
final cones; $W_m=E_R$, $W_0=E_L$, and
$W_j-W_{j-1}=-B_j+H_j$ is supported in $J(y_j)$.

\emph{Locality and descent to polynomial vectors.}
For an adjacent interchange of physical fields in the suffix,
smear all suffix variables with any compact test supported where
that pair is spacelike. Keep this test fixed. The two orderings
have the same partial spectral cones in the prefix and $x$:
$q_h\in\overline V_+$ for $h\le j$ and
$q_j+k\in\overline V_+$. Smearing the suffix does not change
those cones. Their difference is therefore the boundary of a
holomorphic function in the ordered tube of the $j+1$ variables
$y_1,\ldots,y_j,x$, with the last imaginary difference measured
from the real boundary. On the nonempty open set where $x$ is
spacelike to every prefix point, the already proved adjacent
$L$ interchanges move it to the extreme left. At that extreme,
physical locality makes the fixed suffix swap zero, regardless
of the prefix. Boundary uniqueness in this partial tube makes
the difference zero for all prefix and $x$ variables. Since the
fixed suffix test was arbitrary on the spacelike-pair patch,
this proves the required distributional locality. For a swap
in the prefix, hold its test fixed instead and move $L$ to the
extreme right. The common final cones, equivalently the negative
reversed partial sums by total momentum conservation, give the
opposite ordered tube in $x$ and the suffix variables. The same
argument applies.

For a fixed homogeneous list, locality moves $L$ to an extreme
position on a spacelike open set.  Ordered tube uniqueness therefore
shows that the specified extreme values determine the form uniquely.
In particular, the result is multilinear in the homogeneous labels
and is covariant, independently of the auxiliary causal splittings.
For arbitrary original entries write
\[
 A_i=\sum_{\Delta\in S_i}A_{i,\Delta},\qquad |S_i|<\infty,
\]
and take the finite sum of the constructed forms over
$S_1\times\cdots\times S_m$, using the same joint test kernel in
every term.  The common ordered spectral cones, locality, wave
equation, and extreme pairings survive this finite sum.  Its dilation
covariance is the termwise action on the labels, not homogeneity of
one total degree.  The distribution orders and constants may depend
on this fixed finite list.

The vector-dependence assertion uses exactly the preceding choice of
which variables to hold fixed, now for the whole finite summed form.
If a fixed ket polynomial $A\Omega$ is zero, even by cancellation
between component terms or word lengths, vary only the opposite
prefix and $x$.  Once the entire ket has been smeared and held fixed,
the finite sum has the same prefix-and-$x$ tube, regardless of the
word lengths on that fixed side.
On the open set where the $L$ interchanges through that prefix
are spacelike, the extreme-left value pairs $l$ with
$C^*A\Omega=0$. The partial-tube argument just given makes
the insertion zero everywhere. It never requires a null
polynomial vector to remain zero under independent translations
of its individual factors. If the bra polynomial is zero,
hold it fixed and use the extreme right and the opposite tube.
For a permitted joint kernel, distributional continuity first
supplies the summed form and its common partial-tube boundary for
that fixed whole kernel.  The zero extreme pairing is then used for
the whole kernel, not for its approximants.  Hermiticity and the
stated covariance follow from uniqueness of the whole form.
This proves
a statement about forms with one $L$, without a positivity
assertion for repeated insertions.

In particular, the construction may now be differentiated in the
$L$ slot and paired with arbitrary physical polynomial vectors.
These are precisely the operations needed to form the conserved
currents below; none requires composing two $L$ insertions.
\end{proof}

\subsection{Conserved charge forms and their vacuum flux}

Use Proposition~\ref{C:single-insertion} to define the conserved
one-insertion forms
\begin{equation}\label{C:formal-currents}
 t_{\mu\nu}=T_{\mu\nu}
 +\frac13(\partial_\mu\partial_\nu-\eta_{\mu\nu}\Box)L,
 \qquad J_\mu=V_\mu-\partial_\mu L.
\end{equation}
Thus $t$ is symmetric and traceless and $\partial^\mu J_\mu=0$.
For a conformal Killing vector $X$, $X^\nu t_{\mu\nu}$ is conserved.
Its commutator flux through a compact Cauchy slab around a compact
physical polynomial $A$, paired with $C\Omega$, defines a charge
form $q_X(C,A)$. The cutoff equals a time step on the relevant causal
support. Locality and conservation permit compact deformations of
this slab. To move the charge from ket to bra requires an additional
statement: a complete vacuum flux through a physical product must
vanish even when $X$ is quadratic.

The compact commutator flux and this complete vacuum flux have
different roles.  Conservation and locality make the former
independent of a compact deformation.  The latter is a limit over
increasing scales and requires an infrared estimate.  We establish
that estimate uniformly in the mass band so that the bands can
subsequently be removed.

\Needspace{6\baselineskip}
\begin{lemma}[Low-energy estimate]\label{C:low-energy}
Let $Z$ be an even Hermitian compact polynomial vacuum vector, smooth
under translations. For an invariant-mass spectral band
$I\Subset(0,\infty)$ and $\Lambda>0$,
\begin{equation}\label{C:low-energy-bound}
 \|1_{P^0<\Lambda}E_I Z\|
 \le C_Z\Lambda\,\langle Z,P^0E_I Z\rangle^{1/2},
\end{equation}
where the constant is independent of $I$ and $\Lambda$.
\end{lemma}

\begin{proof}
The mass projection commutes with the wedge modular data and hence
preserves all wedge real spaces containing $Z$. Wedge separation
therefore makes the self-commutator function of $E_IZ$ vanish at
bounded times and sufficiently large spatial separation, with the
spatial bound independent of $I$. Its time derivative at zero has
compact spatial support and is bounded in absolute value by
$2\langle Z,P^0E_IZ\rangle$. If $\mu_I$ is the positive translation
spectral measure of $E_IZ$, the spatial Fourier transform of this
derivative, after division by its fixed imaginary factor, is the
sum of the positive marginal measures of $p^0d\mu_I(p)$ at
$\boldsymbol p$ and $-\boldsymbol p$. Fourier transformation of the
compactly supported bounded function bounds each of these positive
measures by
$C_Z\langle Z,P^0E_IZ\rangle d^3\boldsymbol p$.
On the positive-energy cone $p^0\ge|\boldsymbol p|$. Hence
\[
 \int_{p^0<\Lambda}d\mu_I
 \le C_Z\langle Z,P^0E_IZ\rangle
 \int_{|\boldsymbol p|<\Lambda}\frac{d^3\boldsymbol p}{|\boldsymbol p|}
 \le C'_Z\Lambda^2\langle Z,P^0E_IZ\rangle.
\]
The integrable singularity at $\boldsymbol p=0$ causes no problem;
the positive mass band has no atom there. Taking square roots proves
the estimate.
\end{proof}

A fully scaled smooth vacuum flux for a quadratic $X$ is of the
form $R U_D(\log R)F$, where $F$ is a fixed purely timelike,
translation-smooth wave. For a compact mass band of $F$, the paired
part of $Z$ lies in $I_R=[c/R,C/R]$. Equation
\eqref{C:low-energy-bound}, with $\Lambda=N/R$, gives a bound
proportional to
$N\langle Z,P^0E_{I_R}Z\rangle^{1/2}$, which tends to zero.
To remove the compact-energy and compact-mass restrictions, write
$F_N=1_{N<P^0<2N}F$ for dyadic $N>0$. The same argument without
a mass restriction bounds the pairing on this dyad by
$C_ZN\|F_N\|\langle Z,P^0Z\rangle^{1/2}$, uniformly in $R$.
The stress-tensor vacuum amplitude is $O((p^0)^2)$ and the
four-momentum measure of a dyadic cone segment is $O(N^4)$.
Thus $\|F_N\|=O(N^4)$ at zero. At infinity the smooth flux
test gives arbitrarily rapid decrease. Hence
$\sum_N N\|F_N\|<\infty$. On a fixed dyad, compact timelike
mass approximation has arbitrarily small norm, and its complement
is controlled by the same uniform bound. Dominated convergence
now proves the full flux assertion. Lower-degree $X$ have a smaller scaling
factor. Polarization into even Hermitian tests proves that every
complete vacuum flux under consideration vanishes.

The left and right charge prescriptions consequently agree and
$q_X$ is Hermitian. For an affine Killing field the improvement
changes the current by the superpotential
\begin{equation}\label{C:affine-superpotential}
 \frac13\partial^\lambda
 \bigl(X_\lambda\partial_\mu L-X_\mu\partial_\lambda L
             +(\partial_\lambda X_\mu)L\bigr).
\end{equation}
Its compact commutator flux is zero. Thus the translation and
Lorentz charge forms are the physical ones. Integration by parts
for dilatations gives
\begin{equation}\label{C:Dprime}
 D'=D+Q,\qquad Q=q_J.
\end{equation}
At this point $D'$ and $Q$ are charge forms, not asserted
self-adjoint operators. Their dependence on two separately
translated polynomial tests is tempered and has each test's ordinary
vacuum translation spectrum. For the bra spectrum one lets the
compact flux follow the ket, translating the polynomial vector
field $X$ as well; Hermiticity supplies the ket spectrum. In
particular strictly energy-damped collective translations are
legitimate arguments of these forms.

\subsection{Identification of the quadratic charge and its domains}

The comparison with the physical current will first be made on
opposite wedge polynomials, smoothed under translations and rotations.
Compact smooth field tests have both kinds of derivatives; the angular
derivatives will control the distant caps. We keep
the large longitudinal cutoff and the growing transverse cutoff
separate: the first is removed at fixed transverse radius, and only
then is that radius allowed to grow.

Write
\[
 k_n(x)=2(n\cdot x)x-x^2n
       =|y|^2n+2uy+2u^2\bar n.
\]
We shall identify its charge with the already constructed positive
self-adjoint moment
$K_n=b(|Y|^2+\kappa)$. First take opposite compact Hermitian wedge
polynomials, as in the bounded-profile Ward identity of
Section~\ref{sec:null-current}. Locality identifies the commutator
flux through $u<0$ with the extreme vacuum commutator of
$(t\cdot k_n)\Omega$ against their product. Translation smoothing
makes these matrix elements smooth and supported in the union of
the corresponding causal supports. Conservation gives a plane
term $|y|^2t_{vv}$ and the following cap terms.

On the lower cap, the nonzero part with positive step function is to
the future of the right wedge polynomial. It has bounded $u$,
$|y|\asymp R$, and $v\gtrsim R^2$. Product cutoffs can therefore be
chosen with $\partial_u\chi_R=0$ there. On a dyad $v\asymp S$,
$S\gtrsim R^2$, the transverse cap amplitudes obey
\begin{equation}\label{C:cap-amplitude}
 C_NRS\bigl(R^2b\sqrt{ab}+Rab+a\sqrt{ab}\bigr)(1+bS)^{-N}.
\end{equation}
The factors follow respectively from the three terms of $k_n$,
the stress-tensor component bounds, the transverse cutoff
derivative, and the cap volume. Integration in $v$ gives the last
factor. With sufficiently many inverse energy weights these
amplitudes have uniformly bounded Hilbert norms. More explicitly,
$p^0$ is comparable to $a+b$, and the allowed transverse momenta
have area at most $Cab$. In the $d^4p$ realization the squared
norms of the three terms of \eqref{C:cap-amplitude}, weighted
by $(1+p^0)^{-m}$ with $m>5/2$, are bounded respectively by
\begin{equation}\label{C:cap-norms}
 C\bigl(R^6S^{-3},\ R^4S^{-2},\ R^2S^{-1}\bigr).
\end{equation}
For example the first squared integrand, after transverse
integration, is bounded by
$CR^6S^2a^2b^4(1+a+b)^{-2m}(1+bS)^{-2N}$;
its $b$ integral is $O(S^{-5})$ and its $a$ integral is finite.
The other two have powers $a^3b^3$ and $a^4b^2$.
Since $S\gtrsim R^2$, all three bounds are uniform.
There is also a positive power gain when paired with a test vector
smooth under rotations. Fix a sufficiently small $\varepsilon>0$.
The region $a<S^{-\varepsilon}$ gains at least
$S^{-3\varepsilon/2}$ in norm by the same integral;
$b>S^{-1+\varepsilon}$ is negligible by increasing $N$.
In the remaining region
$b/p^0\le C S^{-1+2\varepsilon}$, so the momentum direction lies
in a spherical cap of power-small area. Rotation Sobolev estimates
bound the spectral projection of the fixed test vector onto this
cap by a positive power of its area. This follows by averaging
the squared projected norms over rotations and then applying
Sobolev evaluation on the rotation group; the same estimate holds
after each of the finitely many required rotation derivatives and
energy weights. The angular cap has area
$O(S^{-1+2\varepsilon})$, so its projection norm is
$O(S^{-1/2+\varepsilon})$. For example any
$0<\varepsilon<1/2$ gives
$\gamma=\min(3\varepsilon/2,1/2-\varepsilon)>0$.
The paired dyads are bounded by $CS^{-\gamma}$, and summing
$S\gtrsim R^2$ gives $O(R^{-2\gamma})$.

A regulator at large $v$ can first be removed at fixed $R$.
Its derivative has size $S^{-1}$, its volume is $O_R(S)$, and the
orbital coefficients are then bounded independently of $S$.
The rapid $bS$ decay gives a vanishing energy-weighted norm.
On the plane, locality restricts the $|v|\asymp S$ tails to
$|y|\lesssim\sqrt S$. Their amplitudes are bounded by
\begin{equation}\label{C:plane-amplitude}
 C_NS^3b^2(1+bS)^{-N}.
\end{equation}
Here transverse integration bounds the squared weighted norm by
$CS^6\int a b^5(1+a+b)^{-2m}(1+bS)^{-2N}\,\dd a\dd b$,
which is uniform. Restriction to $a<S^{-\varepsilon}$ gives
$O(S^{-\varepsilon})$ in norm. The preceding angular argument
therefore again gives summable
paired tails, uniformly when the growing quadratic spatial
profile is truncated by a bounded-factor smooth cutoff.
The term $\partial_v^2L$ integrates to zero: the derivative of
the separating step is in a spacelike gap. Moving both derivatives
onto a distant tail cutoff gives the scalar-wave amplitude
$C_NS(1+bS)^{-N}$. Its squared weighted norm is bounded by
$CS^2\int ab(1+a+b)^{-2m}(1+bS)^{-2N}\,\dd a\dd b$,
and has the same low-$a$ and angular gains. The surviving plane integral is
therefore the growing second moment of the bounded-profile Ward
identity. It equals the form of $K_n$.

This identity must be extended on a genuine domain. We record the
domain argument rather than treating an unbounded charge form as
an operator.

Its two inputs are geometric Tomita localization and smoothness under
the boost group.  They first yield power domains for vectors with
strictly one-sided supports.  Positive-energy analyticity then
transports these estimates to damped translates, where the charge
identity can be compared on one common dense domain.

\Needspace{6\baselineskip}
\begin{lemma}[An affine Hardy estimate]\label{C:affine-Hardy}
Let $B=-i\partial_\tau$ on $L^2(\R,\dd\tau;\mathcal K)$, where
$\mathcal K$ is any Hilbert space. Suppose
$f\in\bigcap_m H^m(\R;\mathcal K)$ and
$g(\tau)=e^{-ice^\tau}f(\tau)\in\Dom e^{-\pi B}$, with $c>0$.
For every $r>0$ there is a finite integer $m$ such that
\begin{equation}\label{C:Hardy-power-bound}
 \|e^{r\tau}f\|\le C_{r,c}
 \bigl(\|f\|_{H^m}+\|e^{-\pi B}g\|\bigr).
\end{equation}
The constants are uniform when $c$ varies in a compact subinterval
of $(0,\infty)$.
\end{lemma}

\begin{proof}
Choose $\chi\in C_c^\infty(\R)$ equal to one on $[0,1]$, put
$\chi_j(\tau)=\chi(\tau-j)$, and write $L=e^j$, $j\ge0$.
At Fourier frequencies $\zeta<-\delta L$, split the convolution
$\widehat\chi_j*\widehat g$ at source frequency $\nu=-\delta L/2$.
The Tomita-domain assumption gives
\[
 \|\mathbf1_{\nu<-\delta L/2}\widehat g\|
 \le e^{-\pi\delta L/2}\|e^{-\pi B}g\|.
\]
For the remaining source frequencies, $|\zeta-\nu|>\delta L/2$.
The Schwartz $L^1$ tail of $\widehat\chi_j$ is $O_N(L^{-N})$,
uniformly in $j$, since translation changes only its phase. Young's
inequality therefore gives
\[
 \|\mathbf1_{\zeta<-\delta L}\widehat{\chi_jg}\|
 \le C_NL^{-N}\bigl(\|g\|+\|e^{-\pi B}g\|\bigr).
\]
On $\zeta\ge-\delta L$, take $\delta$ smaller than half the minimum
of $ce^{\tau-j}$ on $\supp\chi_j$. The phase
$-ce^\tau-\zeta\tau$ then has derivative of absolute value at least
$C(L+|\zeta|)$. Repeated integration by parts in the Hilbert-valued
integral gives
\[
 \|\widehat{\chi_jg}(\zeta)\|
 \le C_N(L+|\zeta|)^{-N}\sum_{k\le N}\|f^{(k)}\|_2.
\]
Its squared integral on this half-line is $O(L^{-2N+1})$ times the
displayed Sobolev bound squared. Since $\|\chi_jf\|=\|\chi_jg\|$,
these two estimates give arbitrary inverse powers of $e^j$ for the
localized norm. Choose $N>r+1$ and sum over $j$; on $\tau<0$ the
weight is at most one. This proves \eqref{C:Hardy-power-bound} and
the stated uniformity.
\end{proof}

\Needspace{6\baselineskip}
\begin{lemma}[Domains for the null moment]\label{C:K-domains}
A smooth compact physical polynomial vacuum vector whose supports
lie strictly on one side of a null plane is in $\Dom K_n^r$ for
every $r>0$, locally uniformly in such placements. Every strictly
energy-damped translate of a compact physical polynomial vector
also has this property. Applying $K_n^r$ to the latter vectors
preserves strong tube holomorphy and the indicated one-sided
boundary continuity.
\end{lemma}

\begin{proof}
On the lower side choose $c>0,d\in\R$ so that all supports lie
in $W_{cp_0+d}$. Tomita affiliation, after undoing
$U_H(c)U_b(d)$, gives the domain condition for $e^{-\pi B}$.
For a compact joint kernel, the same condition is the closed
Tomita-graph limit of factorized tests from
Section~\ref{sec:foundations}.  Collective boost derivatives act on
the smooth joint kernel, so their norms are controlled by the old
joint-distribution seminorms.  No affiliated joint-smear operator is
needed for this vector statement.
Remove the fixed constant translation $U_b(d)$ first; it commutes
with $H$ and preserves boost smoothness of these field vectors.
On the positive part of the affine representation use
$H=e^\tau$, $B=-i\partial_\tau$. The resulting vector $f$ satisfies
\[
 e^{-ice^\tau}f\in\Dom e^{-\pi B},
\]
where the original vector $f$ has Sobolev derivatives of every
order in $\tau$. Lemma~\ref{C:affine-Hardy} gives
$e^{r\tau}f\in L^2$ for every $r>0$. The $H=0$ part is
harmless. Since $K_n\le H$ strongly and the two commute, this
proves the assertion for the lower placement. Reflection and
equality of the reflected moments prove the upper placement.
The choices have positive geometric margins, giving locally
uniform bounds.

For $U(x+i\eta)A\Omega$, $\eta\in V_+$, consider the bounded
Hilbert-valued holomorphic function
\[
 1_{(L,\infty)}(K_n)U(x+s\eta)A\Omega,
 \qquad \operatorname{Im}s>0.
\]
A real interval of $s$ places the entire translated support on
one side of the plane. On that interval its norm is $O(L^{-N})$
for every $N$, locally uniformly; elsewhere it is bounded by
the original vector norm. The subharmonic logarithmic norm
estimate with harmonic measure of that interval yields any
prescribed inverse power at every fixed interior point, by
starting with a sufficiently large $N$. Spectral-tail
integration gives all the power domains, locally uniformly.
Closedness of $K_n^r$ and uniform convergence of spectral
truncations give strong holomorphy and the stated boundaries.
\end{proof}

Let $\mathscr D$ be the linear span of all vectors obtained from these
strictly energy-damped compact polynomial translates by finitely many
real translations, dilatations, and compact smooth dilation convolutions
\[
 \int \rho(s)U_D(s)\xi\,\dd s,\qquad
 \rho\in C_c^\infty(\R).
\]
The locally uniform bounds of Lemma~\ref{C:K-domains} and closedness
allow each such integral in the graph of every required $K_n^r$, also
after translating the origin; iteration gives the same conclusion for
every finite sequence of these operations. The space is dense and
invariant under real translations and dilatations by construction.
Its defining smooth kernels also
place it in $\Dom D$; the compact dilation convolutions form a graph
core for $D$, so $\mathscr D$ is such a core.
The preceding charge identity extends to $\mathscr D$ by
tube uniqueness. Thus
\begin{equation}\label{C:charge-identity}
 q_{k_n}(C,A)=\langle C\Omega,K_nA\Omega\rangle
\end{equation}
on this common domain, for every origin and future null normal.
The linear relations among conformal vector fields give vector
relations among the $K_n$ there: test them against the dense
domain and use \eqref{C:charge-identity}. In particular
\begin{equation}\label{C:quadratic-charge-translation}
 K_n(s):=U_a(s)K_nU_a(-s)
       =K_n-2s(D'-B)+2s^2a\quad\hbox{on }\mathscr D.
\end{equation}
This equation realizes $D'$ and $Q=D'-D$ as symmetric vector
operators on $\mathscr D$, by a difference of the indicated
operators. Only each $K_n(s)$ is being used as self-adjoint;
no essential self-adjointness of $Q$ is required.

\subsection{The canonical action on the one-field space}

The charge identity supplies more information than the spectral
comparison alone.  Its translated quadratic dependence first kills
the nonnegative comparison defect in the scalar channels.  The
spectral support statement then upgrades equality of compression
forms to a reducing-subspace identity.  Only after that step do we
treat the remaining helicities.

\Needspace{6\baselineskip}
\begin{theorem}\label{C:canonical-Pi}
The one-field space $\Pi$ reduces each $K_n$. On $\Pi$, $K_n$
equals the null special-conformal generator of the auxiliary
positive conformal representation carried by $l,J,t$.
Moreover $\langle Q\phi,\psi\rangle=0$ for
$\phi\in\mathscr D$, $\psi\in\Pi$.
\end{theorem}

\begin{proof}
We first treat the scalar labels $l,j_z,s_z$. The comparison
\eqref{D:order-zero}, together with $c_V=l+j_z$ and
$c_T=l/3+s_z$, gives order-zero comparison on each of them.
For smooth compact waves with $b,M>0$ compare the compression
of $K_n$ with $K_{0,n}=b(|Y|^2+h_0)$.
Their nonnegative expectation difference, evaluated on
$U_a(-s)\psi$, has zero second derivative in $s$.
Indeed the second-difference identity from
\eqref{C:quadratic-charge-translation} transfers to these waves
by symmetry against $\mathscr D$; the domains follow from the
spectral comparison. Direct differentiation gives the same
second difference for $K_{0,n}$. A nonnegative affine function
on the whole real line is constant.

Testing with separated compact multipliers in $b,k$ now shows
that the positive radial form
$h_0-\kappa_{\rm compr}$ is invariant under multiplication by
$e^{itM^2}$. Its distribution kernel is annihilated by
$M^2-M'^2$. On $M,M'>0$ this forces a diagonal kernel; positivity
makes it a positive measure $\nu(dM)$. It is bounded by the
$h_0$ form. For any compact interval $I\Subset(0,\infty)$,
choose logarithmic cutoffs $f_L=1$ on $I$, vanishing outside
successively larger compact intervals, with
\[
 \int_0^\infty M|f_L'(M)|^2\,dM\longrightarrow0.
\]
For example piecewise linear ramps in $\log M$ have energy
$O(L^{-1})$, and smoothing preserves that bound. Consequently
$\nu(I)\le\int|f_L|^2d\nu\to0$. The compression forms agree.

This is also equality of operators on a reducing copy. The
spectral comparison gives
\[
 \|1_{(s,\infty)}(\kappa)\psi\|
 \le\|1_{(s,\infty)}(h_0)\psi\|.
\]
Equality of the first moments makes the squared tail norms
equal almost everywhere. The support inclusion underlying this
inequality and equality in the projection norm inequality give
equality of the projected vectors. Spectral continuity gives
it at all continuity points, hence equality of all spectral
measures and reduction. Comparing the odd translation
differences with the canonical scalar formula shows that
$Q\mathscr D$ is orthogonal to these labels. Varying the null
direction spans all polarizations of $\Pi$, proving the last
assertion of the Theorem.

For clarity, the projection argument above does not require the two
spectral measures to commute.  Write $\iota$ for the scalar-copy
embedding, $E_s=\mathbf1_{(s,\infty)}(\kappa)$ and
$P_s=\mathbf1_{(s,\infty)}(h_0)$.  Spectral comparison gives
$E_s\iota=E_s\iota P_s$.  Equality of the first moments is equality
of the integrals of $\|E_s\iota\psi\|^2$ and $\|P_s\psi\|^2$;
their pointwise difference is nonnegative, so it vanishes almost
everywhere.  Equality in the contraction estimate for the orthogonal
projection $E_s$ therefore gives $E_s\iota\psi=\iota P_s\psi$.
One first uses a countable dense form core and then spectral
continuity to obtain the operator identity for every threshold.
This proves reduction rather than merely equality of expectations.

We recall explicitly the reference generator on the remaining
polarizations. On transverse helicity of absolute weight $d$
it is
\begin{equation}\label{C:canonical-hankel}
 K_{c,n}=b(|Y|^2+h_d).
\end{equation}
For a rank-$r$ primary of dimension $r+2$, $r=1,2$, contract
$r-d$ slots with $n$ and use
$\partial_vO_i-\partial_iO_v$ in each of the other slots,
discarding joint transverse traces. The wave is a constant
multiple of $b^rM^d$ times its helicity label. Under the null
conformal map of Section~\ref{sec:foundations}, the transverse
Jacobian is $\gamma(e_i+2cy_i n)$. The curl cancels the
$y_i n$ terms; derivatives of those terms contribute only
the discarded traces. The transformed wave thus has the
effective scalar weight $2+d$, giving
\eqref{C:canonical-hankel} by the scalar computation. This
uses only the positive primary two-point kernels and also
acts on their descendant combinations $V,T$. Its origin
translation formula contains $D$, without $Q$.

On smooth compact timelike $\Pi$ waves the actual domains follow
from \eqref{F:spectral-equation}. Put
$X_n=K_n-K_{c,n}$ there. Symmetry-testing
\eqref{C:quadratic-charge-translation} on $\mathscr D$, using
the orthogonality just proved, shows that $X_n$ intertwines
translations and is linear in $n$. The same is true on the
energy-damped compact $V,T$ waves in $\mathscr D$: symmetry
now gives $Q=0$ as a vector on these waves, and their reference
domains follow from the primary and descendant tube action.
The translation formula gives polynomial graph bounds.
Integrating translates against Schwartz functions therefore
commutes $X_n$ through compact timelike momentum multipliers.
Such multiples of the damped $V,T$ waves span the smooth
compact timelike tests of $\Pi$.

It follows that $X_n$ is a measurable fiber coefficient map.
Linearity in $n$ supplies a covector index $\lambda$. The
coefficient tensors
\[
 X_\lambda V_\mu(p)=X_\lambda J_\mu(p),\qquad
 X_\lambda T_{\mu\nu}(p)=X_\lambda t_{\mu\nu}(p)
\]
are locally square integrable up to the nonzero light cone
from inside: this follows first with the damping and smooth
test multipliers, and then locally without them by choosing
those multipliers nonvanishing. Their dimensions are
respectively $d'=2,3$, and they are covariant tensors.
A rest-frame boost component of weight $\omega$ has norm
proportional to $M^{d'-2-\omega}$ when boosted to a fixed
nonzero energy and $M\downarrow0$. The transverse mass
measure is proportional to $M\,dM$. Local integrability
therefore excludes every weight $\omega\ge d'-1$.

For the first tensor, contracting the $\mu$ slot with a null
vector leaves the $\lambda$ covector parallel to its lowered
null vector: its transverse and same-null contractions have
the forbidden weights. At rest conservation makes the
$\mu$ slot purely spatial. Reversing its spatial direction
changes the contracted tensor by a sign, but would require
it to be parallel to a different null covector. It is zero.
For the second tensor, contract both $\mu,\nu$ slots with
the same null vector. The forbidden weights again leave
only that parallel covector. Both slots are spatial at rest,
so reversing direction leaves the contraction unchanged and
again forces it to be zero. Varying the direction and using
symmetry recovers every spatial tensor coefficient. Hence
both tensors vanish and $K_n=K_{c,n}$ on compact timelike
tests.

For $d\ge1$, conjugation by $M^{1/2}$ turns the radial
operator into
$-\partial_M^2+(d^2-1/4)M^{-2}$, which is limit-point at
both endpoints. Compact smooth radial tests, combined with
smooth transverse cores and compact $b$ cutoffs, are thus
graph cores for \eqref{C:canonical-hankel}. The $d=0$ part
was already reducing with its specified Hankel extension.
The closed graph of the actual self-adjoint $K_n$ contains
the graph of the reference self-adjoint operator on each
copy. Resolvent uniqueness gives reduction of their sum
$\Pi$ and equality there, as claimed.
\end{proof}

\subsection{The residual dilation charge vanishes}

The remaining charge $Q$ is not assigned an independent
self-adjoint realization.  Instead, future and past cone comparisons
provide two families of genuine self-adjoint generators whose first
moments contain the same $Q$.  Opposite exponential bounds on their
limiting measures will force both families to concentrate on the
physical dilation spectrum.

Fix $n,\bar n$ and write $K=K_n$. Truncate the future time
cone by $u<R$ and the past cone by $u>-R$. Their actual
real spaces are contained in comparison real spaces
constructed with $K$ in place of the full parabolic
generator. This is \eqref{C:actual-parabolic-localization} for the
translated parabolic cut; reflection treats the past cone.
Let $T_R^+$ be the comparison future modular generator
and let $-T_R^-$ be the comparison past modular generator.
The affine product formulas and
\eqref{C:quadratic-charge-translation} give self-adjoint
operators satisfying
\begin{equation}\label{C:cone-comparison-generators}
 T_R^\pm=D+Q\mp\frac{K}{2R}\quad\hbox{on }\mathscr D.
\end{equation}
For example the future expression is
$B(R)-K(R)/(2R)$; the past modular expression is
$-B(-R)-K(-R)/(2R)$. Their spectral measures have uniformly
bounded second moments on $\mathscr D$ as $R\to\infty$.
The same is true on a smooth single-field vector supported
compactly in either time cone: the relevant planes avoid
its support, so Lemma~\ref{C:K-domains} and the affine
product domains apply. Testing against $\mathscr D$ shows
that
\[
 T_R^\pm\xi\pm K\xi/(2R)
\]
is a Hilbert vector independent of $R$ and of the sign.

\Needspace{6\baselineskip}
\begin{proposition}\label{C:Q-zero}
The charge form $Q$ vanishes on all physical polynomial
vectors. In particular $Q=0$ on $\mathscr D$.
\end{proposition}

\begin{proof}
Average the spectral measures of $T_R^\pm$ over
$S<\log R<2S$, with normalized Lebesgue measure, and take
common weak limit points as $S\to\infty$. One may first
test continuous functions on the compactified real line.
The uniform second moments on dense vectors remove mass
at infinity. The limits are normalized positive
operator-valued measures, not assumed projection-valued.
Dilatations shift $\log R$ by a fixed amount, whose
boundary contribution in the averages is $O(S^{-1})$.
Thus the limits commute with $D$. They have joint
positive measures $E^\pm(dd,ds)$ with the spectral
measure of $D$. Uniform second moments also give uniform
integrability of first moments, including polarized
cross moments on the vectors just described.

Only these first moments will be compared with $D+Q$.  The
logarithmic averaging is used to obtain commutation with the known
self-adjoint operator $D$, while the uniform second moments prevent
loss of first moments in the weak limit.  Neither operation turns
the limiting positive operator-valued measures into spectral measures.

Let $O$ be a Hermitian scaling component of real dimension
$\alpha$ and $\xi=O(f)\Omega$, with $f$ real, smooth and compact
in the future cone. The scaling orbit continues for
$0<\operatorname{Im}z<\pi$ because $e^zx$ has forward
timelike imaginary part there. Compact interior-cone
support and integration by parts give Hilbert boundary
values at both ends; real scale covariance gives uniform
norm bounds along the strip. Consequently
\begin{equation}\label{C:single-scale-strip}
 e^{-\pi D}\xi=e^{i\pi\alpha}O(f^-)\Omega,
 \qquad f^-(x)=f(-x).
\end{equation}
We need this strip identity also for complex linear combinations of
real dilation translates. Write $\xi_s=U_D(s)\xi$ and let
$\xi_s^-$ be its reflected-test wave. Hermitian conjugation of the
two-point function, with both test coordinates reflected, gives
\[
 \langle\xi_s^-,\xi_t^-\rangle=\langle\xi_t,\xi_s\rangle.
\]
For $p(D)=\sum_jc_jU_D(s_j)$ put
$p^\#(D)=\sum_j\overline{c_j}U_D(s_j)$; the real parameters $s_j$
are unchanged. With $\omega_\alpha=e^{i\pi\alpha}$ and
$\eta=e^{-\pi D}\xi=\omega_\alpha\xi^-$, the reflected Gram identity yields
\begin{equation}\label{C:filtered-adjoint-norms}
 \|p(D)e^{-\pi D}\xi\|=\|p^\#(D)\xi\|,
 \qquad \|p^\#(D)\eta\|=\|p(D)\xi\|.
\end{equation}
For all sufficiently large truncations, the comparison Tomita partner
of $p(D)\xi$ is $p^\#(D)\xi$, and that of $p(D)\eta$ on the past
side is $\omega_\alpha^{-2}p^\#(D)\eta$. Their norms are therefore exactly
the two norms in \eqref{C:filtered-adjoint-norms}. The finite set of
dilated supports remains compactly in the respective cone. In
particular, we do not require a complex filter to preserve its real space.

Real-space containment and these Tomita norm bounds, followed by the
spectral-measure limits in the future and past cones, give
\begin{align}
 \int |p(d)|^2e^{-2\pi s}\,\dd\langle E^+\rangle_\xi
 &\le\|p(D)e^{-\pi D}\xi\|^2,\label{C:Jensen-plus}\\
 \int |p(d)|^2e^{-2\pi d}e^{2\pi s}
                  \,\dd\langle E^-\rangle_\xi
 &\le\|p(D)\xi\|^2.\label{C:Jensen-minus}
\end{align}
Truncation followed by monotone convergence justifies
the unbounded exponential tests in the limiting measures.
The first moments in $s$, with every such filter, agree
for the two signs by
\eqref{C:cone-comparison-generators} and the uniform
moment bounds. Exponential polynomials determine finite
measures: polarization recovers every exponential
$e^{itd}$, and Fourier uniqueness then recovers the
measure. Thus these equalities and inequalities hold
conditionally on the common $d$ marginal $\mu_\xi$.

To spell out the first-moment assertion, the finite filters $p(D)$
preserve the class of smooth cone-supported single-field vectors;
they also commute with $e^{-\pi D}$.  Thus both signs admit every
same filter.  Their limiting first-moment measures in $d$ are finite
by the uniform second-moment bounds.  Taking
$p(d)=1+z e^{itd}$ and polarizing in $z\in\mathbb C$ shows that
these two measures have identical Fourier transforms, hence agree.
This is equality of conditional means, rather than just equality
of their integrals over $d$.

Disintegrate the two scalar joint measures as
$\mu_\xi(dd)\nu_d^\pm(ds)$. Their conditional first
moments $m_\pm(d)$ agree almost everywhere, while
\eqref{C:Jensen-plus}--\eqref{C:Jensen-minus} say
\[
 \int e^{-2\pi s}\nu_d^+(ds)\le e^{-2\pi d},
 \qquad
 \int e^{2\pi s}\nu_d^-(ds)\le e^{2\pi d}.
\]
Jensen gives $m_+(d)\ge d\ge m_-(d)$.
Equality of the means forces both to equal $d$; strict
convexity then forces both conditional laws to be
$\delta_d$. This argument requires only positivity and
normalization of the measures, not multiplicativity of
the operator-valued limits. Positivity also gives
Cauchy--Schwarz for each measurable set, so cross
measures against $\xi$ are supported on $s=d$.
Their first moments, paired with $\phi\in\mathscr D$,
now give
the following identity.  Indeed, if a measurable set is disjoint
from $s=d$, its positive operator has zero quadratic value on $\xi$,
and Cauchy--Schwarz makes every mixed value against $\xi$ zero as
well.  The mixed measure is therefore concentrated on the diagonal;
its $d$ marginal is the mixed spectral measure of $D$.  Its first
moment is consequently that of $D$, while the uniformly integrable
first-moment limits from \eqref{C:cone-comparison-generators}
give the form of $D+Q$.  Subtracting the two gives
\begin{equation}\label{C:Q-single-orthogonality}
 \langle Q\phi,\xi\rangle=0.
\end{equation}

The strip and Jensen argument just completed used a Hermitian
homogeneous field of dimension $\alpha$ with a real test. Splitting a
complex test into its real and imaginary parts extends
\eqref{C:Q-single-orthogonality} to all tests. A complex homogeneous
field is the sum of its Hermitian real and imaginary parts,
\[
 O_\Delta=\frac{O_\Delta+O_\Delta^\dagger}{2}
       +i\,\frac{O_\Delta-O_\Delta^\dagger}{2i},
\]
which have the same grade because adjunction preserves real scaling
dimensions.  Linearity therefore extends
\eqref{C:Q-single-orthogonality} to every homogeneous original field.
For $O=\sum_{\Delta\in S_O}O_\Delta$, its finite scaling support then
extends the conclusion to $O$ with the same smearing test.  This
step sums conclusions; it does not assign the nonhomogeneous field
one strip phase or apply Jensen's inequality to cross-grade terms.

It remains to pass from single fields to products.
Define the insertion $\delta_QO(y)$ by the commutator
with $J(y+r)$ integrated against a compact flux cutoff
in $r$. Conservation and locality shrink the flux as
needed. This insertion is relatively local to physical
fields and is tempered with the ordered spectra: in
this cutoff operation the adjacent momenta are first
combined, so the spectator spectral cones are unchanged.
Its extreme vacuum values against every physical
polynomial vanish by
\eqref{C:Q-single-orthogonality} and the charge identity.
To justify that passage, damp both collective
translations first, remove the ket damping on the
smooth future-cone single-field test, and remove the
bra damping distributionally. Translation covariance
removes the future-cone support restriction. Move the
insertion to the extreme on a spacelike open set;
ordered tube uniqueness makes the zero equality global.
Finally the commutator flux obeys the derivation rule
over every finite list of product slots.  This rule is the
telescoping of the adjacent single-insertion commutators just
constructed, not a new operator axiom for $J$ or $L$.  For a joint
test kernel the telescoping identity is first an identity of joint
distributions and is then tested with that same kernel.  Each
summand is zero, proving the
Proposition.
\end{proof}

\subsection{Modularity of time cones and the actual time action}

We now return from comparison real spaces to algebras of actual
bounded operations.  The equality $Q=0$ determines the limiting
comparison generator, but identifying the physical cone modular
group also requires a KMS identity and separation of the vacuum.
Both will be obtained for the original cone algebra before using
half-sided modular inclusion for line diamonds.

\Needspace{6\baselineskip}
\begin{theorem}\label{C:cone-modularity}
The vacuum is cyclic and separating for the actual
future and past cone algebras. Their modular generators,
with the normalization of Section~\ref{sec:foundations},
are respectively $D$ and $-D$.
\end{theorem}

\begin{proof}
By Proposition~\ref{C:Q-zero},
$T_R^\pm\psi\to D\psi$ on $\mathscr D$. This is a
core for $D$: it is dense and invariant under its
unitary group, and smoothing by compactly supported
smooth functions of the group parameter supplies a
graph core. The self-adjoint core convergence criterion
therefore gives $T_R^\pm\to D$ in strong resolvent
sense, and their real unitary groups converge strongly.

Take two bounded operations in a fixed compact
subregion of the future cone. For all sufficiently
large $R$, their vacuum and adjoint-vacuum vectors
belong to the comparison standard space and its
complexification with the required Tomita domains.
Writing the two operations as $A,B$, the standard-space
two-vector identity gives the scalar holomorphic continuation
of $\langle B^*\Omega,e^{itT_R^+}A\Omega\rangle$ to
$0\le\operatorname{Im}z\le2\pi$. Its values are defined by
\[
 F_R(t+iy)=
 \big\langle e^{-yT_R^+/2}B^*\Omega,
 e^{itT_R^+}e^{-yT_R^+/2}A\Omega\big\rangle,
 \qquad 0\le y\le2\pi,
\]
with the opposite boundary
\[
 F_R(t+2\pi i)
 =\langle e^{itT_R^+}A^*\Omega,B\Omega\rangle.
\]
One obtains the width $2\pi$ by assigning up to $\pi$
of imaginary evolution to each vector and using
$e^{-\pi T_R^+}=J_RS_R$; the comparison conjugation
disappears from the displayed boundary formula. The split-vector
pairing uses only $\Dom e^{-\pi T_R^+}$; no vector
$e^{-2\pi T_R^+}A\Omega$ is required.
The bounds depend only on the four vector norms.
Normal-family convergence
on the strip, together with the strong convergence on
the real edge, gives the vacuum KMS identity for $D$
on these actual operations. The identical argument
in the past gives the identity for $-D$. Bounded
strong-star approximation extends the identity from
compactly localized operations to the cone algebras.

Scaling already acts as an automorphism of each actual
cone algebra. The vacuum is cyclic by Reeh--Schlieder.
For clarity, faithfulness is part of the conclusion:
the KMS identity implies that a bounded operation
annihilating the vacuum has adjoint annihilating it.
Apply this to $BA$ for arbitrary $B$ if $A\Omega=0$.
Then $A^*B^*\Omega=0$; cyclicity gives $A^*=0$.
Thus the vacuum is separating. Uniqueness of the
modular automorphism group of this faithful vector
state identifies the generators as claimed.
\end{proof}

\Needspace{6\baselineskip}
\begin{theorem}[Actual time-M\"obius action]\label{C:time-mobius}
For every future unit timelike vector $e$, the actual
Hilbert space carries a positive-energy representation
of the universal covering group of the time-M\"obius
group. Locally it maps physical diamond algebras with
tips on the line $\R e$ to the corresponding diamond
algebras. Its affine subgroup is the physical time
translation and dilation subgroup. The canonical lift of
the connected stabilizer of each finite interval is its
actual diamond modular group.
\end{theorem}

\begin{proof}
Choose two line parameters $a<b$. Let $A,B_t$ be the
dilation generators centered at $ae,be$, and let $C$
be the modular generator of the physical diamond with
those tips. Inclusion of the diamond into the future
cone at $ae$ and the past cone at $be$ is half-sided
modular by Theorem~\ref{C:cone-modularity} and isotony.
The common vacuum is cyclic and separating for both
inclusions; no cyclicity of a relative commutant is
needed. The half-sided inclusion theorem gives the
affine relations corresponding to
\begin{equation}\label{C:three-mobius-fields}
 A=(z-a)\partial_z,\qquad
 B_t=(z-b)\partial_z,\qquad
 C=\frac{(z-a)(b-z)}{b-a}\partial_z.
\end{equation}
More specifically, $A-C\ge0$ scales with weight $-1$
under $A$, with the negative inward inclusion convention
$N=U_{A-C}(-1)MU_{A-C}(1)$; and $-B_t-C\ge0$ scales
with weight $+1$ under $-B_t$. The pair $A,B_t$
already has the physical affine relation with
$P_e\ge0$.

These are relations of unitary groups, not just
commutators on a tentative common domain. Ordered
products of the three classical one-parameter groups
are coordinates near the identity. The three affine
relations give exact pair-switching identities, so
reordering products proves a local unitary
homomorphism. It integrates uniquely to the simply
connected cover. This is the three-subgroup
integration criterion of
\cite[Lemma~1.1]{GuidoLongoWiesbrock1998}; see also
\cite[Lemmas~2.3--2.4]{MorinelliTanimoto2019}.
In its circle formulation the intervals are $(a,b)$, $(b,\infty)$,
and $(\infty,a)$, with stabilizer fields $C,B_t,-A$. Their three
pair relations are precisely the two cone--diamond relations above
and the physical affine relation between the two cone groups.
Its translation subgroup is the physical one, so
the representation has positive energy.

Affine covariance transports $C$ to the modular
generator of every line diamond. A sufficiently small
M\"obius transformation relative to a fixed finite
interval factors into an affine map of its tips and
an element of its connected stabilizer, with the lift
near the identity. Both factors have the
asserted action on its actual algebra. This proves
the local geometric action, without replacing the
physical net by a larger net.

The construction is made for each chosen time direction using its
actual cone and diamond algebras.  Its local geometric statement is
the one needed below: the connected interval stabilizer acts by that
diamond's modular group.  Compatibility sufficient to integrate all
spacetime conformal transformations on the full net has not been
used in this argument.
\end{proof}

We finish by recording the representation data needed for extraction.
Separate the trivial representation. The nontrivial irreducible
positive-energy representations of the simply connected cover are
the lowest-weight modules with $h>0$
\cite[Theorem~1.5.3 and the following direct-integral discussion]{LongoLectures2008}.
The general representation may contain a direct integral of these
modules, rather than a discrete sum. To obtain the decomposition,
disintegrate the central unitary of a full rotation. Within each
central character, compact energy lies in one integer translate;
positivity and the raising and lowering norm relations decompose it
into orthogonal lowest-weight ladders. The real ladder coefficients
give measurable multiplicity spaces $\mathcal K_h$.

Each nontrivial irreducible module restricts to the unique strictly
positive irreducible affine representation
\cite[Corollary~2.33 and equation~(6)]{NeebOlafsson2017}.
Multiplication by $E^{1/2}$ converts its $L^2(\dd E)$ realization
to the common affine realization
\begin{equation}\label{C:time-primary-waves}
 \mathcal H_h=L^2(\R_+,\dd E/E),\qquad
 E=P_e,\qquad
 (U_D(s)\psi)(E)=\psi(e^sE),\qquad
 \psi_{h,z}(E)=E^he^{izE}.
\end{equation}
The last wave, for $\operatorname{Im}z>0$, transforms as a primary of
weight $h$: its inner product kernel is
\[
 \int_0^\infty E^{2h-1}e^{iE(z-\bar w)}\dd E
 =\frac{\Gamma(2h)}{[-i(z-\bar w)]^{2h}},
\]
with the branch fixed in the upper half-plane. Kernel covariance on
the dense span of these waves fixes the primary realization and its
lift to the cover.
Here $h$ denotes a lowest-weight parameter, distinct from the internal
parabolic operator used in the null-cut comparison.
On the nontrivial summand, put
\[
 \mathcal K=\int_{h>0}^{\oplus}\mathcal K_h\,\dd\nu(h),\qquad
 \cH_{\mathrm{nontriv}}=L^2(\R_+,\dd E/E)\otimes\mathcal K,
 \qquad \mathsf h v_h=h v_h.
\]
Thus a measurable weight-band projection is
$\id\otimes\mathbf1_I(\mathsf h)$. No discreteness of this weight
operator follows from, or is required by, the field-label spectrum.

\Needspace{6\baselineskip}
\begin{lemma}\label{C:mobius-real-spaces}
On any bounded weight band contained in $h>0$, the
real standard space of a line diamond is generated
by the primary waves \eqref{C:time-primary-waves}
smeared with real smooth compact interval tests and
multiplicity vectors real for the appropriate phased
conjugation. Weight-band projections commute with
the diamond Tomita operator.
\end{lemma}

\begin{proof}
Let $J$ be the actual future-cone modular conjugation. The
positive-translation relation gives $JP_eJ=P_e$. The half-sided
cone--diamond relation gives the same formula for its positive
special-conformal partner, while $JDJ=-D$. These three generating
groups show that $J$ implements time reflection on the integrated
M\"obius representation.

In each lowest-weight ladder choose the canonical reflection
conjugation, measurably in $h$; its existence and multiplicity form
are given in \cite[Theorem~1.6.3 and its universal-cover extension]{LongoLectures2008}.
The product of this conjugation with $J$ commutes with the full
representation. Inequivalent lowest weights cannot be mixed by that
commutant, so $J$ preserves every measurable weight band. In the
common affine realization it is pointwise energy conjugation together
with an energy-independent multiplicity conjugation $j$, preserving
the spectral projections of $\mathsf h$. This argument distinguishes
weights $h$ and $1-h$ even when their Casimir values coincide.

For a real compact half-line test $f$, the strip continuation gives
\[
 e^{-\pi D}\bigl(E^h\widehat f(E)v\bigr)
 =e^{i\pi h}E^h\widehat f(-E)v.
\]
Hence its Tomita fixed-point condition is
$v=e^{-i\pi h}jv$, the phased real structure in the statement.
These tests are complex cyclic by
Fourier uniqueness. Their real closure is invariant
under dilatations and contained in the half-line
real space. Gaussian smoothing in $D$ gives graph
density for Tomita, hence equality of the two real
spaces. On compact positive weight bands these operations have uniform
bounds, and exhaustion treats a bounded band approaching zero.
The half-sided unitary sends the half-line
space onto the finite-interval space and transports
the tests by primary covariance. All these operations
commute with the weight-band projections. Each such projection
therefore commutes with the diamond conjugation, strongly commutes
with its modular operator, and preserves the domain and action of
its Tomita operator.
\end{proof}
